\chapter{事件锚点与研究导航}
\label{app:event-navigation}

本附录只列出已经在连续叙事首次出现处设置的阅读入口。事件账本仍保留完整的覆盖与证据记录；这里不以生成的逐年条目替代正文，也不为尚未写入叙事的事件制造跳转。

\begin{description}
  \item[\eventref{V03-581-ANNAL-CENTRAL}{隋受禅（581）}] 受禅仪文、关中军政资源与接收江南的开端。
  \item[\eventref{ST-617-LI-YUAN-TAIYUAN}{太原起兵（617）}] 唐朝接收隋制之前的军镇、道路与关中条件。
  \item[\eventref{ST-690-ZHOU-FOUNDATION}{武周建国（690）}] 改号、请愿文本与地方服从不能相互替代。
  \item[\eventref{ST-755-AN-LUSHAN-REBELLION}{安禄山起兵（755）}] 三镇资源集中、洛阳失守与潼关压力的转折。
  \item[\eventref{LT-880-062}{黄巢入长安（880）}] 都城失守、成都行在与区域军府重组的开端。
\end{description}

每一入口旁的材料提示只说明能够确证的年序、行动者和空间条件；兵数、伤亡、居民态度与个人动机仍须回到各章的来源说明和限度。
